\documentclass[11pt,letterpaper]{article}
\usepackage[utf8]{inputenc}
\usepackage[T1]{fontenc}
\usepackage{lmodern}
\usepackage[margin=0.28in]{geometry}
\usepackage{hyperref}
\hypersetup{colorlinks=false, urlcolor=black, pdfborder=0 0 0}

\setlength{\parindent}{0pt}
\setlength{\parskip}{0pt}
\renewcommand{\baselinestretch}{0.82}
\setlength{\lineskip}{0pt}
\setlength{\lineskiplimit}{0pt}

\makeatletter
\renewcommand{\@listi}{%
  \leftmargin=\leftmargini
  \rightmargin=0pt
  \labelsep=4pt
  \labelwidth=16pt
  \itemindent=0pt
  \listparindent=0pt
  \topsep=0pt
  \partopsep=0pt
  \parsep=0pt
  \itemsep=-11pt
}
\let\saveditemize\itemize
\let\savedenditemize\enditemize
\renewcommand{\itemize}{\bgroup\setlength{\baselineskip}{0.65\baselineskip}\saveditemize}
\renewcommand{\enditemize}{\savedenditemize\egroup}
\makeatother

\begin{document}

{\Large \textbf{HSU, Ke-Cheng (Clio)}}\\[1pt]
(852) 6840 9753 \quad $|$ \quad \href{mailto:khsu@connect.ust.hk}{khsu@connect.ust.hk} \quad $|$ \quad \href{https://github.com/cliohsu1997}{github.com/cliohsu1997} \quad $|$ \quad \href{https://cliohsu1997.github.io/}{cliohsu1997.github.io}

\vspace{0.15pt}

{\large \textbf{Professional Summary}}\\[-0.5pt]
PhD candidate in Economics at HKUST specializing in empirical industrial organization and digital economics. Experienced working with large proprietary datasets, including e-commerce platform transactions and retail point-of-sale records. Builds reproducible, testable structural-estimation pipelines on HPC/SLURM and AWS, and communicates results with clear data visualizations. Experience estimating demand and pricing models and delivering counterfactual welfare and policy analysis for platform and retail markets, including identifying firm conduct parameters, merger simulation, and market entry/exit counterfactuals conducted in PhD-level industrial organization coursework. Seeking an economics consulting or research economist role.

\vspace{0.15pt}

{\large \textbf{Education}}\\[-0.5pt]
\textbf{The Hong Kong University of Science and Technology}, Hong Kong\\
\textit{Doctor of Philosophy in Economics} \hfill Expected Graduation: June 2027\\
Research interests: Empirical industrial organization and digital economics.\\
\textit{Master of Philosophy in Economics} \hfill 2020--2022\\
CGA: 4.1/4.3.\\
\textit{Bachelor of Science in Mathematics and Economics, Minor in Actuarial Mathematics} \hfill 2015--2019\\
CGA: 3.9/4.3.\\
Awards: Dean's List (2016--2019); University Scholarship (730K HKD); Academic Achievement Medal; Postgraduate studentship (1.5M HKD); Redbird scholarship (40K HKD).

\vspace{0.15pt}

{\large \textbf{Experience}}\\[-0.5pt]
\textbf{Research Assistant}, National Taiwan University \hfill 2020\\
\vspace{-6pt}
\begin{itemize}
  \item Modeled crisis intervention as regulator--creditor negotiation; derived and simulated equilibrium predictions to support policy design.
\end{itemize}
\textbf{Teaching Assistant \& Lecturer}, Department of Economics, HKUST \hfill 2020--present\\
\vspace{-6pt}
\begin{itemize}
  \item Lecturer for MSc Math Camp (4.53/5); TA for PhD and undergraduate economics (avg 4.72/5).
\end{itemize}

\vspace{0.15pt}

{\large \textbf{Projects}}\\[-0.5pt]
\textbf{Return Policy Design in E-Commerce} \textit{(Job Market Paper)}\\
\vspace{-6pt}
\begin{itemize}
  \item Built and estimated a structural model of purchase, insurance take-up, and return decisions under monopoly pricing, using within-retailer switches between automatic and optional coverage on Taobao (2017--2023; \(\sim\)62M RMB revenue; six apparel categories; \(\sim\)20--30\% returns).
  \item Developed a reproducible, testable estimation pipeline on AWS and HPC/SLURM: version-controlled steps that clean and link orders, prices, insurance, and returns; scraped and OCR-digitized product attributes in Python; accelerated Monte Carlo and constrained optimization in R/C++, with unit-tested code and recovery checks on simulated data.
  \item Estimated category-specific demand, consumer uncertainty about product fit, and return responses using method of moments (GMM); used bootstrap inference to rank policies by seller profit, consumer surplus, and social welfare.
  \item Quantified a purchase--return efficiency trade-off: automatic coverage expands valuable purchases in high-margin categories, while optional coverage outperforms it by \(\sim\)15--30 pp of first-best welfare in thinner-margin categories by curbing excessive returns; partial RPD (discounted price + \(\sim\)10--15\% restocking fee) beats both outside dresses, attaining \(\sim\)75--78\% of first-best welfare.
\end{itemize}
\vspace{-3pt}
\textbf{Inventory and Pricing under Resale Price Maintenance (RPM)} \textit{(w/ K.\ Kawaguchi, HKUST; K.\ Kawai, U.\ Tokyo)}\\
\vspace{-6pt}
\begin{itemize}
  \item Specified a dynamic oligopoly model of pricing and inventory under demand uncertainty and RPM, with equilibrium stockouts and unsold returns (35--40\% of Japanese books printed are discarded), to quantify pro-competitive inventory effects versus waste from price commitment.
  \item Built title-by-market-by-month analysis datasets from monthly bookstore point-of-sale data (deliveries, sales, returns; 1{,}000+ stores) through a tested, end-to-end pipeline: linked products across sources via standardized ISBNs; drew representative samples across book categories; used Hausman-style price instruments to address price endogeneity; collected local demographic data to capture municipality-level market heterogeneity in the estimation; adjusted for launch timing using Amazon release dates.
  \item Tackled endogeneity in demand and supply regressions with instrumental variables (2SLS) and structural model inversion; validated the estimator on simulated data (Monte Carlo); applied LASSO and panel fixed effects to select demographic controls; implemented iterative simulation of equilibrium in C++.
  \item Showed that the classical static model overestimates marginal cost by \(\sim\)20\% when ignoring dynamic demand uncertainty and inventory decisions.
  \item Produced descriptive statistics and short-turnaround figures (sales/return/stockout rates, category demand distributions, substitution patterns) and counterfactuals on deposit-refund policies that cut unsold stock.
\end{itemize}

\vspace{0.15pt}

{\large \textbf{Skills}}\\[-0.5pt]
\textbf{Languages:} Mandarin (Native), English (Fluent), Cantonese (Basic).\\
\textbf{Technical:} Python, R, C++, \LaTeX, AWS, HPC/SLURM; MLE, GMM, minimum distance, EM algorithm, bootstrap.\\
\textbf{References:} Kei Kawai (UTokyo; Compass Lexecon), Kohei Kawaguchi (HKUST; UTokyo Econ Consulting), Pak Hung Au (HKUST).

\end{document}
